% !TeX program = xelatex
% Compile: tectonic django_migrations_report.tex
% Or: xelatex django_migrations_report.tex (twice).
\documentclass[11pt,a4paper]{article}
\usepackage{fontspec}
\setmainfont{Times New Roman}
\setsansfont{Arial}
\setmonofont{Courier New}
\usepackage{polyglossia}
\setdefaultlanguage{russian}
\newfontfamily\cyrillicfonttt{Courier New}
\setotherlanguage{english}
\usepackage[margin=21mm,headheight=15pt]{geometry}
\usepackage{xcolor,fancyhdr,enumitem,booktabs,tabularx,longtable}
\usepackage{fvextra}
\usepackage{titlesec}
\usepackage{needspace}
\usepackage{hyperref}
\definecolor{accent}{HTML}{164B63}
\definecolor{muted}{HTML}{52616B}
\hypersetup{colorlinks=true,linkcolor=accent,urlcolor=accent,citecolor=accent,pdftitle={Миграция данных и версионирование схем БД в Django},pdfauthor={Сценарий для двух спикеров}}
\urlstyle{same}
\setlength{\parindent}{0pt}
\setlength{\parskip}{5pt}
\linespread{1.07}
\setlength{\emergencystretch}{3em}
\pagestyle{fancy}
\fancyhf{}
\fancyhead[L]{\small\sffamily Django: миграции и безопасный деплой}
\fancyhead[R]{\small\sffamily Сценарий · 30 минут}
\fancyfoot[C]{\small\thepage}
\titleformat{\section}{\Large\sffamily\bfseries\color{accent}}{\thesection.}{0.5em}{}
\titlespacing*{\section}{0pt}{8pt}{9pt}
\DefineVerbatimEnvironment{code}{Verbatim}{fontsize=\footnotesize,breaklines=true,breakanywhere=true,frame=single,rulecolor=\color{muted},framesep=5pt,baselinestretch=1}
\newcommand{\speaker}[2]{\Needspace{5\baselineskip}\par\medskip{\sffamily\bfseries\color{accent}Спикер #1\hfill #2}\par\smallskip}
\newcommand{\cue}[1]{{\small\itshape\color{muted}[#1]}\par}
\newcommand{\src}[1]{{\small\color{muted}Источники: #1.}}
\newcommand{\parthead}[2]{\clearpage\section{#1}\textcolor{muted}{\sffamily #2 · две реплики по 75 секунд}\par\medskip}
\begin{document}
\hypersetup{pageanchor=false}
\begin{titlepage}
\vspace*{18mm}
{\sffamily\color{accent}\large СЦЕНАРИЙ ВЫСТУПЛЕНИЯ}\par
\vspace{12mm}
{\sffamily\bfseries\Huge Миграция данных\par и версионирование\par схем БД в Django}\par
\vspace{9mm}
{\sffamily\Large django-migrations, стратегии rollback\par и работа с большими таблицами}\par
\vspace{18mm}
{\large Полный русскоязычный доклад для двух спикеров}\par
\vspace{4mm}
30 минут · 12 частей · 24 чередующиеся реплики\par
\vfill
\textbf{Техническая база:} Django 5.2, PostgreSQL 17; Alembic рассматривается как инструмент экосистемы SQLAlchemy. Это выбранные версии для воспроизводимого объяснения, а не утверждение о самых новых версиях.\par
\textbf{Дата проверки источников:} 12 ноября 2026 года.\par
\textbf{Формат:} основной текст произносится; код показывается и объясняется по смыслу. Ремарки, библиография и приложение вслух не читаются.\par
\vspace{6mm}
{\small Первичные источники: документация разработчиков и публичные отчёты команд об инцидентах. Все примеры кода авторские, адаптированы к единому учебному проекту.}
\end{titlepage}
\hypersetup{pageanchor=true}
\section*{Как пользоваться сценарием}
\addcontentsline{toc}{section}{Как пользоваться сценарием}
Реплики строго чередуются: 1, 2, 1, 2 на протяжении всего доклада. Каждая рассчитана на 1 минуту 15 секунд, включая передачу слова и указанные паузы на код. Целевой темп основного текста: около 130--140 слов в минуту; сложные термины произносятся медленнее. Основной произносимый текст содержит около 3580 слов: при таком темпе это 26--28 минут речи, остальное время отведено показу кода, паузам и передаче слова. Общий тайминг включает объяснение примеров, но не живое выполнение команд и не вопросы из зала. Перед выступлением нужна одна совместная репетиция с таймером.

Сквозной пример: приложение \texttt{orders}, модель \texttt{Order}, таблица \texttt{orders\_order}, числовой первичный ключ \texttt{id}. Мы добавляем поле \texttt{external\_id} типа UUID, затем индекс и ограничение обязательности. Названия миграций иллюстративные: в реальном проекте зависимости берутся из его графа. Отдельный пример \texttt{state} относится к небольшому справочному преобразованию.

\begin{tabularx}{\textwidth}{@{}lXl@{}}
\toprule
Время & Часть & Реплики \\
\midrule
00:00--02:30 & Зачем миграциям отдельная дисциплина & 1, 2 \\
02:30--05:00 & makemigrations, migrate и граф & 1, 2 \\
05:00--07:30 & Исторические модели и миграции данных & 1, 2 \\
07:30--10:00 & Что означает частичная миграция & 1, 2 \\
10:00--12:30 & Совместимость версий и expand--contract & 1, 2 \\
12:30--15:00 & Большие таблицы и пакетное заполнение & 1, 2 \\
15:00--17:30 & Индексы и ограничения PostgreSQL & 1, 2 \\
17:30--20:00 & Rollback, roll forward и восстановление & 1, 2 \\
20:00--22:30 & Фикстуры и начальные данные & 1, 2 \\
22:30--25:00 & Alembic рядом с Django & 1, 2 \\
25:00--27:30 & Два реальных инцидента & 1, 2 \\
27:30--30:00 & Проверка, деплой и завершение & 1, 2 \\
\bottomrule
\end{tabularx}

\textbf{Разделение фактов и рекомендаций.} Ссылки в квадратных скобках ведут в библиографию. Описания API и поведения СУБД основаны на документации. Порядок релизов, критерии остановки и схема реагирования представляют инженерный синтез для учебного примера. Исторические инциденты не выдаются за сбои Django.

\textbf{Безопасная демонстрация.} Команды изменения данных показываются на изолированной учебной БД. Примеры не являются готовым планом для произвольного продакшна: размер пакета, таймауты, права, маршрутизация БД и допустимые блокировки выбираются для конкретной системы.

\parthead{Миграция как изменение работающей системы}{00:00--02:30}
\speaker{1}{00:00--01:15}
Представим интернет-магазин: заказы создаются каждую секунду, а нам нужно добавить внешний идентификатор заказа. На ноутбуке это выглядит как новое поле модели и две команды. Но в рабочей системе одновременно существуют миллионы старых строк, новые запросы, фоновые задачи и несколько запущенных версий приложения. Изменение таблицы становится изменением общего договора между всеми этими участниками.

Сегодня мы разберём, как Django хранит историю схемы, как переносить данные и что делать, если переход прервался. Под словом «миграция» будем различать изменение структуры и преобразование содержимого. Добавление столбца относится к первому, заполнение идентификаторов старых заказов ко второму. Оба действия нужны, но у них разная стоимость и разные способы проверки.

Наш главный вопрос звучит так: может ли сервис корректно работать на каждом промежуточном шаге? Если ответ положительный только для конечной схемы, план ещё не готов. Поэтому мы будем обсуждать не только Python-код, но и порядок релизов, конкурентные записи, наблюдаемость и условия остановки.
\speaker{2}{01:15--02:30}
Начнём с трёх представлений одной системы. Первое --- текущие модели в исходном коде. Второе --- состояние моделей, которое Django восстанавливает из файлов миграций. Третье --- настоящие таблицы и ограничения в базе. Для штатного развития проекта они должны согласовываться. Ручной SQL или неверное использование служебных команд способны нарушить это согласие.

Встроенную систему часто называют django-migrations. В этом докладе речь именно о механизме Django, а не о дополнительном пакете, который надо установить. Файлы миграций хранятся вместе с кодом и описывают переходы между состояниями. Однако сама история переходов не является резервной копией данных.

Это различие особенно важно при удалении. Если миграция убрала столбец с адресами клиентов, обратная операция может вернуть столбец, но не прежние адреса. Поэтому выражение «миграция обратима» ещё не означает «бизнес-данные восстановимы». Дальше мы отдельно рассмотрим технический возврат схемы, возврат версии приложения и восстановление содержимого базы. Их нельзя объединять в одну кнопку.
\src{\cite{migrations,operations}} 

\parthead{Две команды и граф зависимостей}{02:30--05:00}
\speaker{1}{02:30--03:45}
Команда makemigrations строит файлы изменений по сравнению моделей с состоянием из истории миграций. Она не является универсальным сравнением моделей с реальной базой. Поэтому ручное добавление столбца через SQL не превращается автоматически в корректно учтённую историю проекта.

Команда migrate исполняет план перехода. Номер файла помогает ориентироваться, но порядок определяет граф зависимостей. Если миграция заказов зависит от изменения пользователей, Django учитывает эту связь. При параллельной работе две ветки могут создать разные конечные узлы одного приложения. Их нужно согласовать, а не просто переименовать один файл.

На экране наш минимальный цикл: создать миграцию, посмотреть SQL и план, затем применить. Мы читаем сгенерированный файл целиком, особенно удаления, переименования и значения по умолчанию. Автоматическое создание экономит работу, но не объясняет намерение разработчика и не оценивает нагрузку на большую таблицу. Передаю слово: что именно следует проверять до исполнения?
\begin{code}
python manage.py makemigrations orders
python manage.py sqlmigrate orders 0002
python manage.py migrate --plan
python manage.py showmigrations orders
python manage.py migrate
\end{code}
\cue{На команды отвести около 10 секунд; не читать их посимвольно.}
\speaker{2}{03:45--05:00}
Первое --- операции действительно соответствуют задаче? Если вместо переименования получились удаление и создание, данные могут оказаться под угрозой. Второе --- на какой базе выполняется план? В проекте с несколькими подключениями нужны осознанный выбор database и согласованные маршрутизаторы. Третье --- какой набор изменений уедет вместе с релизом?

Файлы генерируются в разработке, проходят проверку и доставляются одинаковыми на тестовый и рабочий контуры. Создавать новые версии истории при каждом старте сервера --- плохая организационная схема: результат перестаёт быть заранее проверенным артефактом релиза.

Мы также различаем план и прогноз длительности. Вывод migrate с флагом plan показывает последовательность, но не измеряет ожидание блокировок, скорость диска или распределение данных. Даже просмотр SQL не заменяет испытание на сопоставимом объёме. В нашей команде ревью миграции поэтому отвечает на два вопроса: правильно ли изменится структура и приемлема ли цена перехода. Если второй ответ отсутствует, зелёный результат на пустой базе ещё не даёт основания запускать изменение в час пик.
\src{\cite{admin,migrations,operations}} 

\parthead{Миграции данных и исторические модели}{05:00--07:30}
\speaker{1}{05:00--06:15}
Теперь нам нужно преобразовать данные. Для этого в Django есть RunPython и RunSQL. В Python-функции мы получаем модель через переданный реестр apps. Это историческая модель на соответствующем шаге графа, а не сегодняшняя версия из models.py. Прямой импорт текущей модели способен сломать установку проекта с нуля через несколько месяцев, когда модель уже изменилась.

В примере небольшая таблица: старое состояние new превращается в pending. Мы явно используем подключение schema editor. Сам по себе RunPython не перенаправляет запросы модели на нужный алиас базы. Историческая модель также не обещает наличие нынешних пользовательских методов: преобразование лучше выразить через доступные поля и запросы.

Обратите внимание: обратной функции здесь нет намеренно. Если pending уже встречался до перехода, по конечному значению невозможно определить происхождение строки. Мы честно объявляем потерю различимости, вместо того чтобы писать формально обратный, но неправильный код. Это пример для небольшой таблицы; пакетную обработку большой покажем отдельно.
\begin{code}
from django.db import migrations

def forwards(apps, schema_editor):
    Order = apps.get_model("orders", "Order")
    alias = schema_editor.connection.alias
    Order.objects.using(alias).filter(
        state="new"
    ).update(state="pending")

class Migration(migrations.Migration):
    dependencies = [("orders", "0002_state")]
    operations = [migrations.RunPython(forwards)]
\end{code}
\cue{Показать три места: apps.get\_model, using, отсутствие reverse\_code; около 12 секунд.}
\speaker{2}{06:15--07:30}
Важная часть подготовки --- сформулировать инвариант. Для нашего преобразования это означает: после окончания нет строк со старым значением, а остальные состояния не изменились. Такой критерий полезнее сообщения «скрипт отработал»: процесс мог завершиться, но затронуть неверный набор строк.

Дальше задаём условия повторного запуска. Здесь фильтр по старому значению не затронет уже преобразованные строки, поэтому повтор не удвоит эффект. Но идемпотентность не равна обратимости: операцию можно безопасно повторять вперёд и всё равно нельзя точно отменить назад. Если точный возврат нужен бизнесу, придётся сохранять исходное значение или происхождение изменения.

Не стоит включать в миграцию отправку писем, обращения к внешнему API или вызов текущей бизнес-логики без отдельного обоснования. База и внешний сервис не образуют одну общую транзакцию; письмо не исчезнет при откате SQL. В нашем учебном плане перенос остаётся локальным преобразованием данных, а внешние действия выполняются отдельно. И наконец, тестируем не только свежую установку, но и переход с предыдущей версии на строках со старыми значениями, пустыми полями и пограничными случаями.
\src{\cite{operations,transactions,querysets}} 

\parthead{Частичные миграции: четыре разных ситуации}{07:30--10:00}
\speaker{1}{07:30--08:45}
Словосочетание «частичная миграция» неоднозначно. Первая ситуация --- мы сознательно хотим остановиться на определённом узле истории. Команда migrate с приложением и номером задаёт целевое состояние, а не указание выполнить произвольный кусочек файла. Зависимости могут затронуть другие приложения, поэтому перед таким действием смотрим полный план.

Вторая ситуация --- исполнение оборвалось. На PostgreSQL обычная атомарная миграция выполняется в транзакции, и при ошибке её транзакционные изменения откатываются. Но если мы отключили atomic, ранее завершённые операции могли сохраниться. На MySQL транзакционной отмены изменений схемы в таком виде нет. Нельзя делать вывод о состоянии таблицы только по последней строке лога.

Третья ситуация --- строка истории и реальная схема разошлись. Флаг fake меняет учёт применения, не исполняя соответствующий SQL. Он полезен при контролируемом согласовании уже существующей схемы, но не исправляет её. Если столбец не создан, пометка «применено» не заставит его появиться. Это особенно опасно при попытке быстро убрать ошибку деплоя.
\begin{code}
python manage.py migrate orders 0003 --plan
python manage.py migrate orders 0003
# --fake changes migration records, not the schema.
\end{code}
\speaker{2}{08:45--10:00}
Четвёртая ситуация --- схема готова, но фоновое заполнение данных ещё продолжается. Django может считать миграции применёнными, а бизнес-переход будет не завершён. Например, новый столбец существует, но у части старых заказов идентификатор пуст. Готовность функции должна проверяться отдельным условием, а не только состоянием таблицы миграций.

После прерывания мы сначала останавливаем автоматические повторные попытки, затем сопоставляем журнал, каталог базы и файлы релиза. Нас интересуют существующие столбцы, ограничения, индексы и уже изменённые данные. Только после этого выбираем повтор, исправляющую миграцию или ручное согласование по проверенному плану.

Представим неатомарный файл из двух операций: добавление индекса прошло, заполнение упало. Повтор всего файла может встретить уже существующий индекс. Поэтому удобнее заранее разделять независимые этапы и задавать условия повторного запуска. Нельзя считать отсутствие записи об успешной миграции доказательством отсутствия её последствий. Отдельно помним про fake-initial: проверка существования начальных таблиц не является полным аудитом соответствия всех типов, ограничений и данных. Наш вывод здесь практический: сначала определить, какой именно вид «частичности» возник, и только потом выбирать команду.
\src{\cite{admin,migrations,writing}} 

\parthead{Совместимость версий: expand--contract}{10:00--12:30}
\speaker{1}{10:00--11:15}
Вернёмся к внешнему идентификатору заказа. Нам нужен переход, при котором старые и новые процессы некоторое время сосуществуют. Разобьём его на расширение, перенос и последующее ужесточение. Этот порядок часто называют expand--contract. Он не является специальной командой Django: это способ организовать несколько изменений.

Сначала добавляем допускающее пустоту поле без немедленного требования уникальности. Старое приложение продолжает работать. Затем выпускаем код, который создаёт идентификатор для новых заказов, а при чтении учитывает старые строки без значения. После этого заполняем прошлые данные. И лишь когда все писатели обновились, а проверка показала отсутствие пропусков и дубликатов, включаем финальные ограничения.

Почему не добавить сразу unique и default? Для существующих строк одноразовое заполнение при добавлении уникального поля способно дать одинаковое значение. Официальная инструкция Django специально разбирает эту ловушку. UUID должен создаваться отдельно для каждой старой записи. Кроме того, правильное преобразование данных ещё не означает, что построение ограничения будет дешёвым: физическую операцию планируем отдельно.
\begin{code}
# 0002: intermediate schema
migrations.AddField(
    model_name="order",
    name="external_id",
    field=models.UUIDField(null=True),
)
# Later application model (after final validation):
external_id = models.UUIDField(default=uuid.uuid4, unique=True)
\end{code}
\cue{Верхняя и нижняя части относятся к разным этапам; показать это явно.}
\speaker{2}{11:15--12:30}
Ключевой риск скрывается между релизами. Если остался старый обработчик очереди, он может продолжать создавать строки без идентификатора. Разовый успешный backfill этого не исправит. Нам нужен барьер: все пути записи переведены на новую логику, старые процессы завершены, и только после этого проводится заключительная проверка.

Для переименования или смены формата задача сложнее. Можно временно хранить старое и новое представление, писать оба и проверять их согласованность. Но двойная запись тоже требует дизайна: ошибка второго обновления не должна оставлять противоречивую пару. В одной базе это часто решается общей транзакцией; при разных хранилищах понадобится другой протокол.

В нашем плане сначала переключаем запись, затем догоняем историю, затем переключаем чтение. Старое представление удаляем отдельным поздним релизом, когда истёк согласованный период возврата. Для каждого этапа формулируем допустимые сочетания версии кода и схемы. Такая матрица сразу обнаруживает опасный сценарий: откатили приложение, а обязательный новый столбец больше не позволяет ему создавать заказы. Значит, быстрый возврат кода нужно проверять ещё до начала миграции, пока схема остаётся расширенной и совместимой.
\src{\cite{writing,fields,alembiccookbook}} 

\parthead{Большие таблицы: управляемое заполнение}{12:30--15:00}
\speaker{1}{12:30--13:45}
Большая таблица --- не только число строк. Для планирования важны размер записи, частота изменений, доступный диск, скорость репликации и допустимая задержка запросов. Один массовый update может создать долгую транзакцию и большой объём журнальных записей. Поэтому переносим данные пакетами и наблюдаем за ценой каждого пакета.

В примере следующая миграция неатомарная, но каждый пакет обёрнут в собственную транзакцию. Сначала выбираем не более тысячи ключей с пустым external id, генерируем разные UUID и обновляем строки. Условие пустоты повторено в update: конкурентный писатель мог уже заполнить значение после нашей выборки. После сбоя повторный запуск снова выбирает незаполненные строки.

Полный код из приложения реализует небольшой пакетный алгоритм, но не обещание достаточной производительности для любой базы. Он делает отдельный update на строку и потому годится как понятная отправная точка, а не как универсальный скоростной загрузчик. Для действительно огромного массива мы вынесем задачу из блокирующего шага деплоя в управляемый фоновый процесс.
\begin{code}
# Core of the batch; full migration is in the appendix.
with transaction.atomic(using=alias):
    for pk in ids:
        qs.filter(pk=pk, external_id__isnull=True).update(
            external_id=uuid.uuid4()
        )
# Select up to 1000 remaining NULL rows before each batch.
# Migration.atomic = False; each batch commits separately.
\end{code}
\cue{На объяснение кода оставить около 20 секунд; обратимость не заявлена.}
\speaker{2}{13:45--15:00}
Почему мы снова выбираем пустые строки, а не используем offset? После каждого пакета набор кандидатов сокращается, и смещение способно пропускать записи. Но у нашего простого решения есть другая цена: повторные поиски оставшихся пустот могут становиться дорогими. Для большого объёма нужны проверенный план запроса и, например, обход диапазонов первичного ключа с сохранением прогресса и заключительной сверкой.

Промежуточный индекс тоже не бесплатен: его создание имеет собственную стоимость. Поэтому решение принимаем по измерению, а не по привычке. Фоновому процессу задаём ограниченный размер работы, возможность приостановки, повторяемость и отчёт о количестве оставшихся строк. Финальная миграция обязана проверить результат, а не просто довериться статусу очереди.

В качестве практического ориентира полезен подход GitLab к пакетным фоновым миграциям: небольшие идемпотентные задания, повтор после ошибки, контроль нагрузки и отдельная проверка завершения. Их реализация относится к Rails, но организационный принцип переносим. Для нашего Django-сервиса заранее выбираем метрики остановки: задержку реплик, задержку запросов и свободный диск. При ухудшении уменьшаем темп или ставим перенос на паузу; молчаливое продолжение ради красивого процента готовности не является целью.
\src{\cite{writing,transactions,querysets,gitlabbatches}} 

\parthead{Индексы, ограничения и блокировки}{15:00--17:30}
\speaker{1}{15:00--16:15}
Теперь заполнение завершилось, и нам нужен индекс. На PostgreSQL обычное создание индекса блокирует запись на время построения. Конкурентное построение позволяет записям продолжаться, но требует больше работы и может ждать старые транзакции. Оно не означает полного отсутствия блокировок или нулевой нагрузки.

Для Django есть AddIndexConcurrently. Его нужно выполнять вне общей транзакции миграции, поэтому в примере atomic равен false. Это отдельная миграция индекса, а не тот же файл, где меняется структура и переносится вся история заказов. Операция учитывает индекс и в состоянии Django, поэтому нам не нужно вручную имитировать это через посторонний SQL.

После прерывания конкурентного построения проверяем каталог: может остаться невалидный индекс. Простое добавление if not exists не доказывает пригодность существующего объекта. Разбираемся в его состоянии и только затем выбираем удаление с повторным построением либо иной документированный способ восстановления. Уникальный индекс требует особенно внимательной проверки дубликатов и конкурентных писателей.
\begin{code}
from django.db import migrations, models
from django.contrib.postgres.operations import AddIndexConcurrently

class Migration(migrations.Migration):
    atomic = False
    dependencies = [("orders", "0003_fill_external_id")]
    operations = [AddIndexConcurrently(
        model_name="order",
        index=models.Index(
            fields=["external_id"], name="order_ext_lookup"
        ),
    )]
\end{code}
\cue{Этот пример создаёт обычный индекс поиска, а не уникальное ограничение.}
\speaker{2}{16:15--17:30}
Для ограничений полезно разделить установку правила и проверку накопленной истории. PostgreSQL допускает not valid для check и внешних ключей: новые изменения проверяются, а старые строки проходят отдельную validation. Это не универсальный параметр для unique или not null. Чтобы сделать поле обязательным, можно сначала проверить соответствующий check, затем запланировать установку not null с учётом блокировок.

Если используем ручной SQL, обязаны согласовать реальную базу с состоянием Django. Для этого существуют RunSQL со state operations и SeparateDatabaseAndState. Но ошибочное описание состояния опасно само по себе. В приложении к докладу приведён минимальный вариант контрольного check через штатные операции Django, а не набор SQL-команд без учёта истории.

Даже короткая операция может долго ждать блокировку. Для сессии миграции задаём обоснованный lock timeout и отдельно statement timeout: один ограничивает ожидание блокировки, другой длительность команды. Во время работы смотрим активные запросы и блокирующие процессы. Числа в учебном примере не норматив: слишком маленький лимит постоянно срывает полезную работу, слишком большой позволяет задержке превратиться в пользовательский сбой.
\src{\cite{pgindex,pgalter,pglocks,pgtimeouts,pgstats,pgoperations}} 

\parthead{Три смысла rollback}{17:30--20:00}
\speaker{1}{17:30--18:45}
Предположим, новая версия приложения ошибается при обработке заказов. Первый вариант --- вернуть предыдущий код, сохранив расширенную схему. Если мы спроектировали совместимость, это часто самый быстрый способ восстановить сервис. Возвращать базу назад только потому, что возвращается приложение, не обязательно.

Второй вариант --- развернуть миграции в обратную сторону до выбранного состояния. Django использует обратные операции, а для RunPython требуется reverse code. Однако возвращение столбца не восстанавливает удалённые значения. И формальная пустая обратная функция, noop, означает именно отсутствие действий: она не доказывает, что преобразование данных отменено.

Третий вариант --- восстановить данные из резервной копии или средствами point-in-time recovery. Для PostgreSQL это требует заранее подготовленного базового бэкапа и непрерывной последовательности архивированных WAL. Такой путь нужно реально испытывать: наличие файла резервной копии ещё не даёт известного времени восстановления. Здесь речь уже о содержимом кластера и выбранной точке времени, а не о простом переключении номера миграции.
\begin{code}
python manage.py migrate orders 0002 --plan
# Execute only after checking reverse operations and dependencies:
python manage.py migrate orders 0002
\end{code}
\speaker{2}{18:45--20:00}
Как выбирать? Сначала определяем цель: восстановить доступность, вернуть правильные значения или устранить несовместимость. Затем оцениваем, какие данные появились после релиза. Если восстановить всю базу на прошлый момент, более поздние корректные записи тоже не окажутся в восстановленном состоянии. Их дальнейшая судьба требует отдельного плана согласования.

Альтернатива обратному переходу --- исправление вперёд, roll forward. Мы выпускаем новую миграцию или совместимый код, который устраняет ошибку, сохраняя текущие данные. Это не всегда лучше отката, но после необратимого преобразования часто честнее, чем изображать возвращение к прежнему миру.

Для нашего магазина до релиза фиксируем две границы. Пока старое поле сохранено, возможен возврат кода по проверенному пути. После удаления источника данных точный возврат зависит от сохранённого происхождения или резервного восстановления. В план инцидента входят ответственный, команда остановки переноса, проверка целостности и ожидаемое время решения. Такой документ полезен только если он соответствует реальному состоянию системы. Поэтому репетируем его на тестовом контуре, включая сбой посередине пакета и повторный запуск, а не ограничиваемся счастливым сценарием завершения.
\src{\cite{operations,pgpitr,alembiccookbook}} 

\parthead{Фикстуры: данные, но не история переходов}{20:00--22:30}
\speaker{1}{20:00--21:15}
Фикстура --- сериализованный набор объектов: например, JSON со справочником стран или демонстрационными заказами. Команда dumpdata выгружает объекты, loaddata загружает их. Это удобно для тестов и воспроизводимого наполнения учебной базы, но не заменяет граф изменения схемы.

Важное отличие проявляется при повторном выполнении. Загрузка фикстуры может вернуть значения объектов к содержимому файла и затронуть уже изменённые записи. Она не описывает семантику перехода «преобразовать только старое состояние и сохранить пользовательские правки». Именно поэтому для обязательных начальных данных Django предлагает миграцию данных, которая включается в обычный процесс migrate.

Есть и техническая особенность: при десериализации пользовательский save не вызывается так, как при обычном сохранении, а обработчики сигналов получают признак raw. Значит, нельзя молча предполагать, что вся бизнес-логика заполнит недостающие зависимости. Мы заранее определяем, какие данные являются тестовыми, какие принадлежат приложению, а какие редактирует пользователь. Для последних статический файл особенно легко становится источником неожиданной перезаписи.
\begin{code}
python manage.py dumpdata orders.Order --indent 2 > orders.json
python manage.py loaddata orders.json
# Use only against a deliberately chosen database.
\end{code}
\speaker{2}{21:15--22:30}
Представим справочник способов доставки. Для локальной демонстрации фиксированные записи подходят прекрасно. Но в рабочем магазине администратор уже изменил описание и цену доставки. Если новый релиз бездумно загрузит прежнюю фикстуру, он может отменить эти изменения. Здесь задача не в выборе формата JSON, а в правилах владения данными.

Когда начальная запись действительно обязательна, миграция может найти её по устойчивому коду и создать при отсутствии. Если требуется обновление существующей записи, в коде явно выражаем, какие поля разрешено менять. Обратное удаление тоже не должно уничтожать объект, который после создания начал использоваться в заказах.

При переносе данных между окружениями учитываем связи и идентификаторы. Natural keys позволяют выражать некоторые связи через естественные ключи вместо совпадения числовых id, но требуют подходящей модели и сериализатора. Ни фикстура, ни natural key сами по себе не превращают перенос в полный бэкап базы. В нашем проекте правила просты: демонстрационные объекты идут в фикстуры, обязательное изменение справочника --- в проверяемую миграцию, а большие пользовательские наборы --- в отдельный контролируемый процесс переноса с собственными проверками результата.
\src{\cite{fixtures,initial,serialization}} 

\parthead{Alembic и Django: границы ответственности}{22:30--25:00}
\speaker{1}{22:30--23:45}
Часто спрашивают: можно ли вместо встроенных миграций подключить Alembic? Сначала уточним модель работы. Alembic относится к экосистеме SQLAlchemy. Его autogenerate сравнивает метаданные таблиц с базой и формирует кандидат на ревизию. В Django штатный процесс опирается на историю состояний моделей. Это разные источники истины и разные механизмы учёта.

Имена команд похожи по назначению: revision создаёт ревизию, upgrade применяет переход, downgrade выполняет обратный путь. Но наличие функции downgrade не доказывает восстановимость потерянных значений. Аналогично Django, автоматически созданный кандидат необходимо читать и исправлять. В частности, переименования таблиц и столбцов Alembic не распознаёт как полноценное намерение: они могут выглядеть как добавление и удаление.

Для обычного Django-проекта практический выбор --- сохранять встроенный механизм владельцем его схемы. Alembic имеет смысл рядом, если отдельный сервис действительно использует SQLAlchemy и владеет своим набором таблиц. Само присутствие двух библиотек в одном проекте не создаёт автоматического согласования их историй.
\begin{code}
# SQLAlchemy service, not commands for a Django migration graph:
alembic revision --autogenerate -m "add external id"
alembic upgrade head
alembic downgrade -1
\end{code}
\speaker{2}{23:45--25:00}
Представим общую базу двух сервисов. Django обслуживает заказы, а аналитический сервис на SQLAlchemy хранит собственные агрегаты. Сосуществование возможно, если явно разделены владельцы таблиц, права на изменение схемы и зависимости релизов. Но позволять обоим инструментам независимо менять orders order --- значит создавать два несогласованных журнала одной реальности.

Это архитектурный вывод, а не запрет конкретного API. Если команда всё же передаёт таблицу от одного механизма другому, она должна отдельно спроектировать базовое состояние, сверку схемы, порядок переключения и дальнейшее владение. Просто отметить одинаковые номера применёнными недостаточно: инструменты не разделяют общий граф зависимостей.

В рекомендациях Alembic по данным есть полезная для нас идея: отдельно расширить схему, выполнить перенос, затем применить финальные ограничения. Для онлайн-перехода описывается сосуществование старого и нового представлений с усложнением приложения. Мы уже использовали тот же принцип в Django. Следовательно, выбор инструмента решает вопрос хранения истории и генерации операций, но не отменяет инженерной задачи совместимости. Ни один генератор не знает, какой заказ нельзя потерять и сколько минут задержки допустимо пользователю. Эти требования формулирует команда, а затем превращает их в проверяемый порядок действий.
\src{\cite{alembicauto,alembiccookbook}} 

\parthead{Как миграции ломали продакшн}{25:00--27:30}
\speaker{1}{25:00--26:15}
Первый документированный случай --- GitHub, 27 ноября 2021 года. По отчёту самой компании, деградация продолжалась два часа пятьдесят минут. Во время миграции большой MySQL-таблицы на финальном переименовании значительная часть реплик чтения попала в semaphore deadlock. Реплики переходили в crash recovery, нагрузка переносилась на оставшиеся, и возник каскадный эффект. Пострадали основные сервисы GitHub.

Это не инцидент Django и не доказательство, что любая миграция вызывает ту же ошибку. Здесь ценен другой факт: длительная предварительная работа не сделала заключительный шаг безопасным автоматически. Переименование, переключение и поведение реплик входят в миграцию наравне с копированием данных.

Наш инженерный вывод для магазина: испытывать нужно весь переход, включая финальное переключение, а мониторинг должен видеть не только основной сервер. Если реплики не справляются, зелёный статус процесса переноса не означает работоспособность сервиса. Заранее нужен способ остановить дополнительную нагрузку и понятное решение, что делать с переключением. Важен не только момент запуска, но и момент, когда ответственность переходит на новую структуру.
\speaker{2}{26:15--27:30}
Второй случай --- публичный разбор GitLab от 9 июня 2020 года. Post-deployment migration, обновлявшая маршруты для lost-and-found группы и осиротевших проектов, завершилась ошибкой нарушения индекса. В отчёте зафиксирован сбой продакшн-деплоя; при первичном разборе само состояние рабочей среды сочли безопасным. Поэтому мы не будем называть этот эпизод подтверждённым массовым простоем пользователей.

Этот пример показывает другую грань: опасность может быть в существующих данных и правилах уникальности, а не в длительности SQL. Отложенная миграция не становится корректной только потому, что выполняется после выкладки приложения. Для нашего справочника это означает предварительный поиск конфликтов и явную политику их разрешения.

Сравним случаи. В GitHub проблема проявилась на уровне работы СУБД и доступности реплик при завершении перехода. В GitLab зафиксирован отказ преобразования из-за ограничения данных и блокировка процесса выкладки. Общая рекомендация не сводится к «проверяйте внимательнее». Нужны конкретные проверки: ограничения на реальных распределениях, последний шаг переключения, наблюдение за зависимыми компонентами и отрепетированный выход из незавершённого состояния. Масштаб последствий у этих случаев разный, и это различие важно сохранять при обсуждении.
\src{\cite{githubincident,gitlabincident}} 

\parthead{Практический порядок деплоя и выводы}{27:30--30:00}
\speaker{1}{27:30--28:45}
Соберём наш план в последовательность, которую можно проверить. До релиза читаем файлы миграций и SQL, прогоняем установку с нуля и обновление с предыдущей версии. Отдельно проверяем данные с дубликатами, пустыми значениями и уже частично заполненным новым полем. Тестовая база должна отражать характер данных и нагрузку достаточно, чтобы выявить рискованные планы запросов.

В процессе выкладки один назначенный исполнитель управляет изменением схемы. Мы не рассчитываем, что каждый одновременно стартующий экземпляр приложения безопасно запустит тот же процесс. Сначала расширяем схему, затем выпускаем совместимый код, переводим всех писателей и запускаем контролируемый перенос истории.

После заполнения сверяем результат: нет пропусков, нет дубликатов, новые записи соответствуют правилу. Лишь затем вводим обязательность и уникальность, предварительно оценив физические операции. Старые поля удаляем отдельным релизом после окна возврата. Это наш авторский план для примера, а не готовый рецепт на любой случай. Его качество определяется тем, можем ли мы назвать состояние системы и доступное действие после каждого шага, включая любой сбой посередине.
\speaker{2}{28:45--30:00}
Во время выполнения держим рядом четыре ответа. Что сейчас меняется? Как понять, что результат правильный? По какому признаку остановиться? Каким способом восстановить работу? Для нашей команды сигналами служат задержка запросов, ожидания блокировок, состояние реплик, свободный диск и прогресс заполнения. Пороговые значения выбираем по требованиям своего сервиса, а не копируем из чужой презентации.

Если что-то пошло не так, сначала определяем текущее состояние и прекращаем дополнительный ущерб. Возврат кода, обратная миграция, исправление вперёд и восстановление резервной копии --- разные инструменты. Выбор зависит от совместимости и того, какие значения уже нельзя восстановить из текущих данных.

На этом завершим. В Django makemigrations описывает изменение, migrate проводит по графу, а корректность перехода обеспечивает подготовленный нами план. Большие таблицы требуют управляемой работы и наблюдения. Миграции данных требуют исторических моделей и честного отношения к обратимости. Фикстуры помогают переносить объекты, Alembic решает задачи своей экосистемы, но ни один из этих механизмов не заменяет проверку промежуточных состояний. Если на каждом этапе сервис остаётся корректным, а у команды есть проверенный следующий шаг, миграция становится предсказуемой инженерной операцией. Спасибо за внимание.
\src{\cite{admin,pgstats,gitlabbatches}} 

\clearpage
\section*{Приложение: технические примеры для обсуждения}
\addcontentsline{toc}{section}{Приложение: технические примеры}
\textbf{Не входит в 30 минут.} Это справочные фрагменты к основному сценарию. Их не нужно читать со сцены. Полноценные миграционные классы в тексте самодостаточны при наличии описанных моделей и предыдущих узлов; остальные фрагменты прямо обозначены как этапы или команды.

\subsection*{1. Полная миграция пакетного заполнения}
Перед запуском старые писатели должны быть остановлены или обновлены. Иначе пустые значения могут появляться снова, и окончание цикла не докажет завершение перехода. В модели на этом шаге поле допускает NULL. Обратный переход намеренно запрещён: идентификаторы могли уже стать внешними ссылками. Генерация UUID не заменяет заключительное уникальное ограничение.
\begin{code}
import uuid
from django.db import migrations, transaction

def fill(apps, schema_editor):
    Order = apps.get_model("orders", "Order")
    alias = schema_editor.connection.alias
    qs = Order.objects.using(alias)
    while True:
        ids = list(qs.filter(external_id__isnull=True)
                   .order_by("pk")
                   .values_list("pk", flat=True)[:1000])
        if not ids:
            break
        with transaction.atomic(using=alias):
            for pk in ids:
                qs.filter(pk=pk, external_id__isnull=True).update(
                    external_id=uuid.uuid4()
                )

class Migration(migrations.Migration):
    atomic = False
    dependencies = [("orders", "0002_external_id")]
    operations = [migrations.RunPython(fill)]
\end{code}
\clearpage
\subsection*{2. Добавление и отдельная проверка CHECK}
В Django 5.2 PostgreSQL-операции позволяют учесть ограничение в состоянии моделей и разделить две фазы. В реальной модели в \texttt{Meta.constraints} должно быть объявлено такое же ограничение. Начальная миграция \texttt{0005} и следующая \texttt{0006} --- разные файлы: установка правила фиксируется до проверки старых строк. Использование \texttt{condition}, а не устаревающего \texttt{check}, соответствует выбранной версии.
\begin{code}
# orders/migrations/0005_external_id_check.py
from django.db import migrations, models
from django.contrib.postgres.operations import AddConstraintNotValid

class Migration(migrations.Migration):
    dependencies = [("orders", "0004_external_id_index")]
    operations = [AddConstraintNotValid(
        model_name="order",
        constraint=models.CheckConstraint(
            condition=models.Q(external_id__isnull=False),
            name="order_ext_not_null",
        ),
    )]
\end{code}
\begin{code}
# orders/migrations/0006_validate_external_id.py
from django.db import migrations
from django.contrib.postgres.operations import ValidateConstraint

class Migration(migrations.Migration):
    dependencies = [("orders", "0005_external_id_check")]
    operations = [ValidateConstraint(
        model_name="order", name="order_ext_not_null"
    )]
\end{code}
Проверенный CHECK ещё не меняет \texttt{null=True} на \texttt{null=False} в состоянии поля. Финальный \texttt{AlterField} планируется отдельно. PostgreSQL 17 может избежать повторного сканирования при \texttt{SET NOT NULL}, если подходящий валидный CHECK доказывает отсутствие NULL, но блокировку для самой смены метаданных всё равно нужно учитывать. Последующее удаление вспомогательного CHECK также делается отдельной операцией.\cite{pgoperations,pgalter}

\clearpage
\subsection*{3. Уникальность: отдельный физический план}
Простой индекс \texttt{order\_ext\_lookup} из доклада не обеспечивает уникальность. Для большого массива PostgreSQL допускает сначала создать уникальный B-tree индекс конкурентно, затем присоединить его как UNIQUE constraint. Ниже именно SQL-эскиз для непартиционированной таблицы, а не готовая Django-миграция:
\begin{code}
CREATE UNIQUE INDEX CONCURRENTLY order_ext_unique_idx
    ON orders_order (external_id);
ALTER TABLE orders_order
    ADD CONSTRAINT order_ext_unique
    UNIQUE USING INDEX order_ext_unique_idx;
\end{code}
Два шага нельзя бездумно поместить в одну атомарную миграцию: конкурентный индекс запрещён внутри транзакционного блока. После ручного SQL состояние Django согласуется через \texttt{SeparateDatabaseAndState} или \texttt{state\_operations}, с выбранным представлением в модели: например, \texttt{UniqueConstraint}. Не следует одновременно создавать один UNIQUE через SQL и второй автоматически через \texttt{unique=True}. Обратный переход учитывает, что индекс после присоединения принадлежит ограничению. Обычный индекс поиска, если он стал избыточным, удаляется отдельным проверенным шагом.\cite{pgindex,pgalter,operations}

\subsection*{4. Наблюдение и лимиты для сессии миграции}
Пример лимитов ниже иллюстрирует назначение параметров, а не рекомендуемые значения для всех систем. Команды \texttt{SET} действуют на текущее подключение: выполнение их в другом окне \texttt{psql} не настраивает подключение \texttt{manage.py migrate}. Для транзакционной операции допустим \texttt{SET LOCAL} в её транзакции; для конкурентного индекса нужен иной способ настройки той же сессии вне транзакционного блока. После работы сессионные настройки возвращаются в согласованное состояние.
\begin{code}
SET lock_timeout = '3s';
SET statement_timeout = '20min';

SELECT pid, state, wait_event_type, wait_event,
       pg_blocking_pids(pid) AS blocked_by,
       now() - query_start AS query_age,
       left(query, 100) AS query
FROM pg_stat_activity
WHERE datname = current_database()
  AND pid <> pg_backend_pid();
\end{code}
Для просмотра чужих запросов могут требоваться дополнительные права. Обнаружение блокирующего процесса не является разрешением автоматически завершить его: сначала выясняются назначение и последствия остановки.\cite{pgtimeouts,pgstats,pgfunctions}

\clearpage
\subsection*{5. Минимальная проверка данных и истории}
\begin{code}
SELECT count(*) AS missing
FROM orders_order WHERE external_id IS NULL;

SELECT external_id, count(*)
FROM orders_order WHERE external_id IS NOT NULL
GROUP BY external_id HAVING count(*) > 1 LIMIT 10;

SELECT indexrelid::regclass AS index_name, indisvalid, indisready
FROM pg_index WHERE indrelid = 'orders_order'::regclass;
\end{code}
Полные проверки на большой таблице сами потребляют ресурсы: время и способ их запуска тоже планируются. Поиск дубликатов перед ограничением --- предварительная диагностика; от гонки с новыми записями в конечном счёте защищает успешно установленное ограничение.\cite{pgcatalog,pgindex}

\textbf{Проверки перед выпуском.} Статическая проверка истории: \texttt{makemigrations --check --dry-run}. Она выявляет изменения моделей без миграций, но не доказывает соответствие реальной схемы. Проверка поведения: новая установка; переход со старой схемы; прерывание и повтор backfill; совместная работа старого и нового кода; реальный план восстановления. Для необратимого преобразования ожидаемая ошибка reverse тоже является осознанным свойством, а не поводом подставлять фиктивный noop.\cite{admin,operations}

\clearpage
\renewcommand{\refname}{Библиография и первичные источники}
\begin{thebibliography}{99}
\small
\bibitem{migrations} Django Software Foundation. \textit{Django 5.2: Migrations}. История моделей, граф, транзакции, особенности СУБД, squashing.\par\url{https://docs.djangoproject.com/en/5.2/topics/migrations/}
\bibitem{admin} Django Software Foundation. \textit{Django 5.2: django-admin and manage.py}. Флаги команд, целевое состояние, fake и fake-initial.\par\url{https://docs.djangoproject.com/en/5.2/ref/django-admin/}
\bibitem{operations} Django Software Foundation. \textit{Django 5.2: Migration Operations}. RunPython, RunSQL, reverse\_code, SeparateDatabaseAndState.\par\url{https://docs.djangoproject.com/en/5.2/ref/migration-operations/}
\bibitem{writing} Django Software Foundation. \textit{Django 5.2: How to create database migrations}. Уникальные поля, гонки, неатомарные миграции и пакеты.\par\url{https://docs.djangoproject.com/en/5.2/howto/writing-migrations/}
\bibitem{transactions} Django Software Foundation. \textit{Django 5.2: Database transactions}. atomic, транзакционные границы и стоимость длинных транзакций.\par\url{https://docs.djangoproject.com/en/5.2/topics/db/transactions/}
\bibitem{querysets} Django Software Foundation. \textit{Django 5.2: QuerySet API reference}. update, bulk\_update и особенности сохранения.\par\url{https://docs.djangoproject.com/en/5.2/ref/models/querysets/}
\bibitem{fields} Django Software Foundation. \textit{Django 5.2: Model field reference}. UUIDField и default.\par\url{https://docs.djangoproject.com/en/5.2/ref/models/fields/}
\bibitem{pgoperations} Django Software Foundation. \textit{Django 5.2: PostgreSQL database migration operations}. AddIndexConcurrently, AddConstraintNotValid, ValidateConstraint.\par\url{https://docs.djangoproject.com/en/5.2/ref/contrib/postgres/operations/}
\bibitem{fixtures} Django Software Foundation. \textit{Django 5.2: Fixtures}. Загрузка, raw-сигналы и обработка объектов.\par\url{https://docs.djangoproject.com/en/5.2/topics/db/fixtures/}
\bibitem{initial} Django Software Foundation. \textit{Django 5.2: How to provide initial data for models}. Миграции данных и повторная загрузка фикстур.\par\url{https://docs.djangoproject.com/en/5.2/howto/initial-data/}
\bibitem{serialization} Django Software Foundation. \textit{Django 5.2: Serializing Django objects}. Natural keys и сериализация.\par\url{https://docs.djangoproject.com/en/5.2/topics/serialization/}
\bibitem{pgindex} PostgreSQL Global Development Group. \textit{PostgreSQL 17: CREATE INDEX}. Конкурентное построение, ожидания и невалидные индексы.\par\url{https://www.postgresql.org/docs/17/sql-createindex.html}
\bibitem{pgalter} PostgreSQL Global Development Group. \textit{PostgreSQL 17: ALTER TABLE}. NOT VALID, VALIDATE, USING INDEX и SET NOT NULL.\par\url{https://www.postgresql.org/docs/17/sql-altertable.html}
\bibitem{pglocks} PostgreSQL Global Development Group. \textit{PostgreSQL 17: Explicit Locking}. Режимы и конфликты блокировок.\par\url{https://www.postgresql.org/docs/17/explicit-locking.html}
\bibitem{pgtimeouts} PostgreSQL Global Development Group. \textit{PostgreSQL 17: Client Connection Defaults}. lock\_timeout и statement\_timeout.\par\url{https://www.postgresql.org/docs/17/runtime-config-client.html}
\bibitem{pgstats} PostgreSQL Global Development Group. \textit{PostgreSQL 17: The Cumulative Statistics System}. pg\_stat\_activity и ожидания.\par\url{https://www.postgresql.org/docs/17/monitoring-stats.html}
\bibitem{pgfunctions} PostgreSQL Global Development Group. \textit{PostgreSQL 17: System Information Functions and Operators}. pg\_blocking\_pids.\par\url{https://www.postgresql.org/docs/17/functions-info.html}
\bibitem{pgcatalog} PostgreSQL Global Development Group. \textit{PostgreSQL 17: pg\_index}. indisvalid и indisready.\par\url{https://www.postgresql.org/docs/17/catalog-pg-index.html}
\bibitem{pgpitr} PostgreSQL Global Development Group. \textit{PostgreSQL 17: Continuous Archiving and Point-in-Time Recovery}. Базовая копия, WAL и восстановление.\par\url{https://www.postgresql.org/docs/17/continuous-archiving.html}
\bibitem{alembicauto} SQLAlchemy / Alembic authors. \textit{Auto Generating Migrations}. Метаданные и ограничения autogenerate.\par\url{https://alembic.sqlalchemy.org/en/latest/autogenerate.html}
\bibitem{alembiccookbook} SQLAlchemy / Alembic authors. \textit{Cookbook: Data Migrations --- General Techniques}. Разделение схемы и данных, онлайн-перенос.\par\url{https://alembic.sqlalchemy.org/en/latest/cookbook.html#data-migrations-general-techniques}
\bibitem{gitlabbatches} GitLab. \textit{Batched background migrations}. Пакеты, идемпотентность, повтор и контроль завершения. Это источник подхода из Rails, а не Django API.\par\url{https://docs.gitlab.com/development/database/batched_background_migrations/}
\bibitem{githubincident} Scott Sanders / GitHub. \textit{GitHub Availability Report: November 2021}. Опубликовано 1 декабря 2021; инцидент 27 ноября.\par\url{https://github.blog/news-insights/company-news/github-availability-report-november-2021/}
\bibitem{gitlabincident} GitLab Infrastructure Team. \textit{2020-06-09: post-deployment migration failure}, incident issue 2251. Инцидент 9 июня 2020.\par\url{https://gitlab.com/gitlab-com/gl-infra/production/-/issues/2251}
\end{thebibliography}
\small Все ссылки проверены 12.11.2026. Для Django и PostgreSQL использованы фиксированные ветки документации; страницы Alembic и GitLab могут обновляться. Описания инцидентов ограничены опубликованными утверждениями команд. Рекомендации по организации нашего деплоя являются авторским инженерным синтезом, а не цитатой из отчётов.
\end{document}
